\section{强化学习在真实世界中的应用}

本节介绍深度强化学习的三个重要应用：自动驾驶中的端到端 RL、围棋中的 AlphaGo 与 AlphaZero，以及近年来 LLM 对齐中的 RLHF。

\subsection{自动驾驶：从仿真到真实部署}

\begin{figure}[H]
\centering
\includegraphics[width=0.95\textwidth]{slides-images/slide-48.jpg}
\caption{VISTA：MIT 开发的高保真自动驾驶仿真器\protect\footnotemark}
\end{figure}
\footnotetext{Slides 第48页。Amini+ IEEE RA-L 2020.}

Amini 教授介绍了他在 MIT 的研究工作------使用 Policy Gradient 方法训练自动驾驶系统。这项工作的关键创新包括：

\textbf{VISTA 仿真器}：一个基于真实驾驶数据的高保真仿真环境。与传统的计算机图形学仿真不同，VISTA 使用真实驾驶视频来合成新的视角，确保仿真画面具有照片级真实感。

\textbf{端到端 RL 训练}：
\begin{itemize}
    \item Agent 的状态是驾驶摄像头的画面
    \item 动作是连续的方向盘转角
    \item 奖励设计极为稀疏：只在车辆驶出车道或发生碰撞时给予惩罚
    \item 策略网络输出高斯分布的参数，从中采样得到具体的转向角度
\end{itemize}

\begin{figure}[H]
\centering
\includegraphics[width=0.95\textwidth]{slides-images/slide-49.jpg}
\caption{全球首辆完全通过仿真中的 RL 训练并直接部署到真实道路上的自动驾驶汽车\protect\footnotemark}
\end{figure}
\footnotetext{Slides 第49页。Amini+ IEEE RA-L 2020.}

\begin{knowledgebox}{从仿真到现实的迁移（Sim-to-Real Transfer）}
这项工作的突破性在于：模型\textbf{完全在仿真环境中训练}，从未接触真实道路数据，但可以\textbf{直接部署到真实汽车}上并安全行驶。这是全球首辆完全通过仿真 RL 训练并在现实世界中部署的全尺寸自动驾驶汽车。关键在于 VISTA 仿真器的高保真度弥合了仿真与现实之间的差距（Domain Gap）。
\end{knowledgebox}

\begin{figure}[H]
\centering
\includegraphics[width=0.95\textwidth]{slides-images/slide-47.jpg}
\caption{训练 RL 在真实世界中应用的难点：运行策略（步骤2）可能导致危险\protect\footnotemark}
\end{figure}
\footnotetext{Slides 第47页。}

\begin{warningbox}{真实环境中训练 RL 的风险}
Policy Gradient 训练算法的第2步"运行策略直到终止"在真实环境中可能非常危险。对于自动驾驶来说，"终止"可能意味着碰撞或事故。因此\textbf{仿真环境}是 RL 在安全关键领域落地的前提------Agent 可以在仿真中安全地犯错和学习，而不会造成真实的损害。
\end{warningbox}

\subsection{围棋：AlphaGo 与 AlphaZero}

围棋（Go）是 Deep RL 最引人注目的应用之一。围棋的复杂度远超国际象棋------19x19 棋盘上合法位置数约为 $2.08 \times 10^{170}$，比宇宙中的原子数量还多。

\begin{figure}[H]
\centering
\includegraphics[width=0.95\textwidth]{slides-images/slide-52.jpg}
\caption{围棋的搜索空间：合法局面数超过宇宙中原子的数量\protect\footnotemark}
\end{figure}
\footnotetext{Slides 第52页。Source: Wikipedia.}

\textbf{AlphaGo}（Silver+ Science 2016）采用了多阶段训练流程：

\begin{figure}[H]
\centering
\includegraphics[width=0.95\textwidth]{slides-images/slide-54.jpg}
\caption{AlphaGo 的训练流程：从人类专家数据到自我对弈\protect\footnotemark}
\end{figure}
\footnotetext{Slides 第54页。Silver+ Science 2016.}

\begin{enumerate}
    \item \textbf{监督学习预训练}：使用人类专家棋谱，以分类任务训练 Policy Network（预测人类棋手的落子位置）
    \item \textbf{强化学习自我对弈}：用 RL 的 Policy Gradient 方法让 Policy Network 通过自我对弈进一步提升，达到\textbf{超越人类}的水平
    \item \textbf{Value Network 训练}：用自我对弈产生的数据训练 Value Network，评估局面的胜率（回归任务）
\end{enumerate}

AlphaGo 在 2016 年以 4:1 击败围棋世界冠军 Lee Sedol，成为 AI 历史上的里程碑事件。

\textbf{AlphaZero}（Silver+ Science 2018）更进一步，\textbf{完全抛弃了人类数据}，从零开始仅通过自我对弈的强化学习来训练：

\begin{figure}[H]
\centering
\includegraphics[width=0.85\textwidth]{slides-images/slide-57.jpg}
\caption{AlphaZero：完全通过自我对弈从零学习，Elo 评分随训练步数快速上升\protect\footnotemark}
\end{figure}
\footnotetext{Slides 第57页。Silver+ Science 2018.}

\begin{importantbox}{从 AlphaGo 到 AlphaZero 的飞跃}
\begin{itemize}
    \item \textbf{AlphaGo}：需要人类专家棋谱作为起点（监督学习预训练 + RL 微调）
    \item \textbf{AlphaZero}：完全从零开始（tabula rasa），不使用任何人类数据，仅通过\textbf{自我对弈 + 强化学习}达到超越所有前任版本的水平
    \item AlphaZero 证明了：在完美信息博弈中，纯 RL 的自我对弈可以\textbf{发现人类从未想到的策略}，突破人类认知的局限
\end{itemize}
\end{importantbox}

\subsection{本章小结}

Deep RL 在自动驾驶和围棋等领域取得了突破性成果。自动驾驶中的端到端 RL 展示了 Policy Gradient 处理连续控制任务的能力，而仿真器是确保安全训练的关键。AlphaGo 和 AlphaZero 系列证明了 Deep RL + 自我对弈可以在极其复杂的决策问题上达到甚至超越人类水平。


